\documentclass[paper=a4, fontsize=11pt]{scrartcl} % A4 paper and 11pt font size

\newcommand{\assignment}{Final Project Template}

\newcommand{\duedate}{November, 2025}

\input{include/hw-template.tex}

\author{
    \textbf{Karan Kendre, Ritik Mahyavanshi, Atharva Dhumal, Akshay Prajapati} \\ 
    \href{https://github.com/your-repo/Code-Circular-Dependency-Detection}{PROJECT GIT REPOSITORY}
}% INFORMATION

\begin{document}

\maketitle % Print the title

\section{Executive Summary and Abstract}

Circular Dependency Detective is a comprehensive tool that systematically detects circular dependencies in Python codebases using Tarjan's Strongly Connected Components algorithm and provides natural language explanations with refactoring recommendations using Retrieval-Augmented Generation (RAG) enhanced Large Language Models (LLMs). The system employs AST-based code analysis for structure-aware parsing, a dual knowledge base architecture combining code embeddings with refactoring patterns, and local LLM deployment via Ollama. Our evaluation demonstrates 100\% cycle detection accuracy across test fixtures, with sub-millisecond detection performance (0.10--0.24ms) and fast RAG retrieval (4.5--4.8ms). The tool successfully identifies all circular dependencies, correctly classifies severity levels, and provides actionable explanations that help developers understand and resolve dependency issues efficiently.

\section{Motivation and Impact}

\subsection{RAG LLM Projects / PAL Agents}

Circular dependencies are a pervasive problem in software development that lead to import errors, tight coupling, and maintenance difficulties. Traditional static analysis tools can detect these cycles but fail to provide context-aware explanations or actionable refactoring recommendations tailored to specific codebases. Our system addresses this gap by combining graph algorithms with RAG-enhanced LLMs to not only detect all circular dependencies with guaranteed completeness but also explain why they occur and how to fix them.

The domain we selected is Python codebase analysis, focusing on import dependency relationships. Our objective is to build an LLM-powered system that detects all circular dependencies with 100\% completeness using Tarjan's algorithm, provides natural language explanations of cycle root causes with specific file and line references, recommends specific refactoring strategies based on code context and retrieved patterns, prioritizes cycles by severity and impact on codebase health, generates interactive visualizations for dependency analysis, and operates entirely locally using Ollama for privacy and enterprise compatibility.

The impact of this work extends to software engineering teams who struggle with circular dependencies in large codebases. By providing both detection and explanation capabilities, our tool reduces the time developers spend debugging import errors and improves code maintainability. The system's local deployment capability (via Ollama) ensures privacy and allows use in enterprise environments where code cannot be sent to external APIs. Our benchmark results show the system can analyze projects in milliseconds, making it suitable for integration into CI/CD pipelines and IDE extensions.

\section{Background and Related Work}

Our implementation leverages several key technologies and design decisions:

\textbf{Graph Algorithms}: We use NetworkX's implementation of Tarjan's Strongly Connected Components (SCC) algorithm for cycle detection. Tarjan's algorithm guarantees finding all cycles in $O(V + E)$ time complexity, ensuring completeness in cycle detection. This is superior to naive cycle-finding approaches that may miss complex multi-file cycles.

\textbf{AST-Based Analysis}: Python's Abstract Syntax Tree (AST) module enables structure-aware code parsing. Unlike regex-based approaches, AST parsing preserves semantic information, allowing us to distinguish between TYPE\_CHECKING imports, local (function-level) imports, and module-level imports. This context is crucial for accurate severity assessment.

\textbf{Retrieval-Augmented Generation}: Our RAG system employs a dual knowledge base approach with separate collections for code chunks and refactoring patterns, enabling precise retrieval of relevant context. This ensures LLM explanations are grounded in actual code and established refactoring practices.

\textbf{Embedding Models}: We use sentence-transformers with the \texttt{all-MiniLM-L6-v2} model, which provides 384-dimensional embeddings optimized for semantic similarity search. This model balances accuracy and inference speed for local deployment, achieving 4.5--4.8ms retrieval times in our benchmarks.

\textbf{Local LLM Deployment}: Ollama enables local execution of large language models without requiring cloud API access. We use CodeLlama 13B, which is specifically trained on code and provides high-quality explanations for software engineering tasks while maintaining data privacy.

Our data consists of two primary sources:

\textbf{Code Data}: The system analyzes Python codebases provided by users. The code is chunked using AST-based structure-aware chunking, preserving semantic units (classes, methods, functions, import blocks). Each chunk includes source code with line numbers, docstrings and function signatures, import relationships, context headers for semantic search, and parent-child relationships (classes contain methods). The code data is interesting because it requires understanding both syntactic structure (AST) and semantic relationships (imports, dependencies). The structured nature of Python code enables effective chunking that preserves context while maintaining searchable units. Our benchmarks show we create 3--6 chunks per file on average, preserving semantic integrity.

\textbf{Refactoring Patterns}: We maintain a knowledge base of refactoring patterns stored as YAML files in \texttt{knowledge\_base/patterns/}. Each pattern includes problem description and symptoms, solution description with before/after code examples, step-by-step refactoring guides, keywords for semantic matching, and complexity ratings and framework support. The patterns are based on established software engineering practices including Extract Interface pattern (breaking bidirectional dependencies), Lazy Import pattern (deferring imports to function scope), Type Checking Guard pattern (using TYPE\_CHECKING for type-only imports), and Dependency Injection pattern (decoupling through injection).

\textbf{Test Fixtures}: For evaluation, we use test fixtures in \texttt{tests/fixtures/} containing simple\_cycle (2-file bidirectional cycle), complex\_cycle (3-file cycle), and type\_checking\_cycle (cycle using TYPE\_CHECKING guards). These fixtures are included in the repository and can be used for validation and benchmarking.

\textbf{Data Access}: The code data is provided by users analyzing their own projects. The refactoring patterns are included in the repository at \texttt{knowledge\_base/patterns/}. Users can extend the pattern library by adding YAML files following the schema defined in \texttt{src/knowledge/schemas.py}. Test fixtures are available in \texttt{tests/fixtures/} for validation. The complete repository is available at: \href{https://github.com/your-repo/Code-Circular-Dependency-Detection}{Repository Link}

\section{Modeling Methodology}

This section details our algorithms and implementation approach. The system follows a modular architecture with distinct components for parsing, graph construction, cycle detection, scoring, RAG retrieval, and explanation generation.

\textbf{Graph Construction and Cycle Detection}: The \texttt{DependencyGraphBuilder} processes all Python files using AST parsing to extract imports, resolves them to file paths, and constructs a NetworkX directed graph. The \texttt{CycleAnalyzer} uses Tarjan's SCC algorithm to find all cycles, computes metrics (length, coupling density, centrality), and assigns severity levels. Our implementation handles edge cases including TYPE\_CHECKING imports, local imports, and simple 2-node cycles.

\textbf{Code Chunking}: The \texttt{StructureAwareChunker} creates semantic code units preserving AST hierarchy. It extracts docstrings, signatures, and context headers, creating separate chunks for import blocks, classes, methods, functions, and top-level code. Our benchmarks show the chunker creates 3--6 chunks per file on average.

\textbf{RAG System}: Our dual knowledge base maintains separate collections for code chunks and refactoring patterns. For a cycle query, we generate query embeddings, retrieve top-$k$ code chunks filtered by cycle files, retrieve top-$k$ refactoring patterns by semantic similarity, and combine into \texttt{RAGContext} for LLM prompts. Our benchmarks demonstrate consistent RAG retrieval performance of 4.5--4.8ms per query.

\textbf{LLM Explanation Generation}: The \texttt{CycleExplainer} builds prompts containing cycle details, edge information, retrieved code context, and refactoring patterns. We use Ollama with CodeLlama 13B, configured with temperature 0.3 and top-p 0.9 for focused, deterministic output.

\textbf{Severity Scoring}: The \texttt{SeverityScorer} computes weighted scores using four factors: length score (30\%), coupling score (20\%), centrality score (30\%), and import type score (20\%). Severity levels are assigned as LOW (score $< 25$), MEDIUM (25--50), HIGH (50--75), or CRITICAL ($\geq 75$).

\textbf{Reproducibility}: The system is fully containerized with Dockerfile, docker-compose.yml, and automated setup scripts (\texttt{start.sh}, \texttt{run.sh}). The \texttt{start.sh} script pulls Ollama Docker image, downloads CodeLlama 13B model automatically, and starts all required services. The complete implementation and README are available at: \href{https://github.com/your-repo/Code-Circular-Dependency-Detection}{Repository Link}

\section{Data and Data Analysis}

\subsection{Data Source(s)}

Our system operates on two types of data:

\textbf{1. Python Codebases}: Users provide their own Python projects for analysis. The system processes all \texttt{*.py} files recursively. For evaluation, we use test fixtures included in the repository at \texttt{tests/fixtures/}:
\begin{itemize}
    \item \texttt{simple\_cycle/}: 2-file bidirectional cycle (module\_a.py $\leftrightarrow$ module\_b.py)
    \item \texttt{complex\_cycle/}: 3-file cycle (a.py $\rightarrow$ b.py $\rightarrow$ c.py $\rightarrow$ a.py)
    \item \texttt{type\_checking\_cycle/}: Cycle using TYPE\_CHECKING guards (models.py $\leftrightarrow$ services.py)
\end{itemize}

\textbf{2. Refactoring Pattern Knowledge Base}: Curated patterns stored as YAML files in \texttt{knowledge\_base/patterns/}:
\begin{itemize}
    \item \texttt{extract\_interface.yaml}: Breaking bidirectional dependencies with Protocol classes
    \item \texttt{lazy\_import.yaml}: Deferring imports to function scope
    \item \texttt{type\_checking\_guard.yaml}: Using TYPE\_CHECKING for type-only imports
    \item \texttt{dependency\_injection.yaml}: Decoupling through dependency injection
\end{itemize}

\textbf{Citation}: The refactoring patterns are based on established software engineering practices from Martin Fowler's ``Refactoring'' patterns, Python-specific best practices from PEP 484 (Type Hints), and common solutions documented in Python community forums. The codebase itself is open source and available at the repository link. Users can analyze any Python project, making the dataset scale effectively unlimited. Our benchmarks used 7 files from test fixtures, but the system has been tested on projects with 42+ files.

\subsection{Data Analysis and Exploration}

\textbf{Code Structure Analysis}: Before processing, we analyze file count and distribution across directories, import statement frequency and types (absolute, relative, TYPE\_CHECKING, local), average file size and complexity, and import graph density (edges/nodes ratio).

\textbf{Cycle Characteristics}: Our analysis of test fixtures reveals:
\begin{itemize}
    \item \textbf{Cycle Length Distribution}: Test fixtures include 2-file and 3-file cycles. The system correctly identifies both types with 100\% accuracy.
    \item \textbf{Severity Distribution}: TYPE\_CHECKING-only cycles are correctly classified as LOW severity, while runtime cycles receive appropriate severity based on length and centrality. The complex\_cycle fixture shows coupling density of 0.5 (3 edges out of 6 possible), indicating moderate coupling.
    \item \textbf{Centrality Correlation}: Files with high centrality scores are correctly identified as more critical.
\end{itemize}

\textbf{Chunking Analysis}: Our benchmarks show simple files produce 3--4 chunks per file (imports, functions, classes), complex files produce 6 chunks per file (TYPE\_CHECKING blocks, multiple classes, methods), with an average of 4 chunks per file across all fixtures. Total: 28 chunks created from 7 files in test fixtures.

\textbf{Pattern Matching}: The RAG system's pattern retrieval successfully matches lazy import patterns for simple 2-file cycles, extract interface patterns for bidirectional dependencies, and type checking guard patterns for TYPE\_CHECKING cycles.

\section{Results and Evaluation}

We evaluated our system on multiple dimensions using test fixtures containing known circular dependencies. Each fixture subdirectory was analyzed separately to ensure proper import resolution and cycle detection.

\subsection{Evaluation Metrics}

\textbf{1. Cycle Detection Completeness}: 
\begin{itemize}
    \item \textbf{Metric}: Percentage of cycles found vs. ground truth
    \item \textbf{Result}: 100\% - All test cases correctly identified
    \item \textbf{Test Cases}:
    \begin{itemize}
        \item Simple Cycle: 2-file bidirectional cycle - $\checkmark$ Correctly detected
        \item Complex Cycle: 3-file cycle - $\checkmark$ Correctly detected  
        \item Type Checking Cycle: 2-file cycle with TYPE\_CHECKING guards - $\checkmark$ Correctly detected
        \item No Cycle: Acyclic project - $\checkmark$ Correctly identified as cycle-free
    \end{itemize}
\end{itemize}

\textbf{2. Performance Benchmarks}:
\begin{itemize}
    \item \textbf{Graph Building}: 0.9--1.2ms mean across fixtures ($\sim$0.43ms per file)
    \item \textbf{Cycle Detection}: 0.10--0.24ms mean ($\sim$0.15ms average)
    \item \textbf{Code Chunking}: 0.30--1.3ms for 2--3 files, creating 7--12 chunks
    \item \textbf{RAG Retrieval}: 4.5--4.8ms mean per query (consistent across fixture sizes)
\end{itemize}

\textbf{3. Severity Classification}: The system correctly classified all test cases: simple 2-file cycle as LOW, complex 3-file cycle as CRITICAL (coupling density 0.5, centrality 0.5), and type checking cycle as LOW (TYPE\_CHECKING and local imports reduce severity).

\subsection{Subjective Examples}

\textbf{Example 1 - Simple Bidirectional Cycle}: Correctly identified 2-file cycle between \texttt{module\_a.py} and \texttt{module\_b.py} with LOW severity. Detected in 0.11ms, graph built in 0.9ms. Created 7 chunks (3.5 per file) in 0.30ms. Accuracy: 100\% - Matched expected result.

\textbf{Example 2 - Complex Multi-File Cycle}: Found 3-file cycle (a.py $\rightarrow$ b.py $\rightarrow$ c.py $\rightarrow$ a.py) with CRITICAL severity. Detected in 0.24ms, graph built in 1.2ms. Created 9 chunks (3.0 per file) in 0.36ms. Accuracy: 100\% - Matched expected result.

\textbf{Example 3 - TYPE\_CHECKING Cycle}: Identified 2-file cycle between \texttt{models.py} and \texttt{services.py} with LOW severity (correctly classified - TYPE\_CHECKING guard and local import reduce runtime impact). Detected in 0.10ms, graph built in 1.0ms. Created 12 chunks (6.0 per file) in 1.3ms. Correctly identified TYPE\_CHECKING flag and local import flag.

\subsection{Confidence Assessment}

We have high confidence in:
\begin{itemize}
    \item \textbf{Cycle Detection}: 100\% confidence - Algorithm guarantees completeness and test results confirm perfect accuracy
    \item \textbf{Performance}: High confidence - Sub-millisecond detection times consistently measured across all test cases
    \item \textbf{RAG Retrieval}: High confidence - Consistent 4--5ms retrieval times measured across multiple iterations
    \item \textbf{Severity Classification}: High confidence - Correctly classified all test cases based on cycle length and import types
\end{itemize}

\subsection{Aggregated Results}

Across all test fixtures: 7 Python files analyzed, 7 import edges, 3 cycles detected (one per fixture type), 28 semantic code chunks created, 100\% overall accuracy (3/3 test cases passed). Average performance: graph building $\sim$1.0ms per fixture, cycle detection $\sim$0.15ms per fixture, code chunking $\sim$0.65ms per fixture, RAG retrieval $\sim$4.7ms per query.

\section{Conclusions}

This project successfully demonstrates the integration of graph algorithms, AST analysis, RAG systems, and LLMs for solving a real-world software engineering problem. Key learnings include: Tarjan's algorithm provides guaranteed cycle detection completeness, achieving 100\% accuracy in our tests; AST-based chunking preserves semantic structure better than line-based approaches, creating 3--6 meaningful chunks per file; dual knowledge base RAG enables precise retrieval of both code context and refactoring patterns, with 4.5ms average retrieval time; local LLM deployment (Ollama) provides privacy and control while maintaining explanation quality; severity scoring based on multiple factors correctly classifies cycle importance.

The system can be applied in several ways: integration into CI/CD pipelines for automated dependency health checks, IDE extensions for real-time cycle detection and explanations, code review tools to flag circular dependencies before merge, educational tools for teaching software architecture principles, and automated code quality assessment in development workflows.

If we had more time, we would focus on: implementing automated code fixes (currently only provides recommendations), expanding the refactoring pattern library with more strategies, building a web dashboard for interactive cycle exploration, adding support for detecting indirect cycles through multiple import levels, integrating with popular Python frameworks (Django, FastAPI) for framework-specific patterns, and performance optimization for very large codebases (1000+ files). Future directions also include multi-language support (JavaScript, Java, Go), machine learning models to predict cycle severity and optimal refactoring strategies, IDE integration for real-time analysis, and historical analysis to track cycle evolution over time in version control.

\section{Submission Guidelines}

Your in-class delivery comprises 40\% of your project grade and the technical artifacts, demonstration, and real-time capabilities comprise 60\%. More details about the presentation breakdown can be found on the \href{https://course.ccs.neu.edu/cs6120s25/presentation-and-project/}{website}. 

\begin{description}
    \item[40\% Documentation] Written and oral delivery of project rationale, justification, results, and conclusions. 

    \begin{itemize}
        \item 10\% Oral Delivery (\href{https://docs.google.com/presentation/d/1THS63CCLEqvzfadkNAwDTP2hCg9bLNnt9LdAo7dUWNc/edit?slide=id.p1#slide=id.p1}{Deck})
        \item 30\% Written Summary (\href{https://www.overleaf.com/project/67d700c6739786050017acaa}{PDF})
    \end{itemize}

    \item[60\% Project Writeup PDF]  Demonstration, Code, and Corner Case Handling

    \begin{itemize}
        \item \textbf{Code and Data Provenance} There are two components that speak to the transparency of your work: verifiability and reproducibility. 

        \begin{itemize}
            \item \textbf{Verifiability} - Our RAG system provides citations for all LLM outputs through specific file paths and line numbers in explanations (e.g., ``module\_a.py line 2 imports module\_b.py''), retrieved code chunks with source attribution, refactoring patterns with before/after code examples, and verification steps included in every explanation. The LLM is instructed to only reference actual code from the codebase and provided patterns, preventing hallucinations. We flag TYPE\_CHECKING-only cycles as low-confidence to prevent false alarms. Our benchmark results demonstrate 100\% accuracy in cycle detection.

            \item \textbf{Problem Complexity} - The system handles codebases ranging from small test fixtures (7 files) to larger projects (42+ files tested). Test fixtures include simple 2-file cycles, complex multi-file cycles, and edge cases like TYPE\_CHECKING guards. The graph algorithms scale efficiently with O(V+E) complexity.

            \item \textbf{Reproducibility} - Full reproducibility is ensured through Docker containerization with automated setup, complete dependency specification in \texttt{requirements.txt}, automated model download in \texttt{start.sh}, clear README with step-by-step instructions, test fixtures for validation (\texttt{tests/fixtures/}), and benchmark script (\texttt{benchmark.py}) for performance validation. Anyone can reproduce our results by running \texttt{./start.sh} and \texttt{./run.sh analyze <project>} or \texttt{python benchmark.py} for benchmarks.
        \end{itemize}

        These perspectives are intimately tied with the quality of your project. 

        \item \textbf{Accessibility and Performance} The system is accessible through command-line interface for local use (\texttt{python -m src.cli analyze <project>}), Docker deployment for easy setup, GCP deployment scripts for cloud hosting, and interactive HTML visualizations. Performance: Graph building and cycle detection complete in milliseconds even for large codebases. Our benchmarks show 0.9--1.2ms for graph building and 0.10--0.24ms for cycle detection on test fixtures. LLM explanations take 10--30 seconds depending on hardware (Ollama generation time).

        \item \textbf{Accuracy} Our LLM application addresses corner cases: TYPE\_CHECKING-only cycles (flagged as low severity, no runtime impact) - $\checkmark$ Correctly handled; local imports (function-level imports reduce cycle severity) - $\checkmark$ Correctly identified; relative imports (correctly resolved to absolute paths) - $\checkmark$ Working; package imports (\texttt{\_\_init\_\_.py} handling) - $\checkmark$ Supported; external library filtering (standard library and third-party imports excluded) - $\checkmark$ Implemented; simple 2-file cycles - $\checkmark$ 100\% accuracy; complex multi-file cycles - $\checkmark$ 100\% accuracy.
    \end{itemize}
\end{description}

Upload your PDF file, GCP (or other) Endpoint Link, and Github URL to \href{https://www.gradescope.com}{Gradescope} as a \emph{single} submission with all members of the team. On the top of the document, include the teammates and project repository. 

\end{document}



